\chapter{类、接口与值语义：设计量化组件}

\begin{introduction}
  \item 从不变量和用例出发设计类，而不是从字段出发
  \item 用值语义表达事件、订单与成交
  \item 把撮合与组合状态隔离为可替换组件
  \item 避免贫血模型、万能管理器和继承滥用
\end{introduction}

\chaptermeta{理解类、组合、接口和前四章的值语义。}{Python 对象默认通过引用共享；本章优先让领域事实成为独立值，并用类维护复合不变量。}{若调用者必须按特定顺序调用多个 setter，说明边界过浅；从不变量重新设计单一操作。}

\section{类的价值是守住边界}

类不应只是把公开字段换成机械 getter/setter。一个有用的类拥有清晰责任、维护不变量，并通过小型公开接口隐藏内部表示。组合模块的责任是：成交发生后，现金与持仓必须一起改变；调用者不应分别调用 \texttt{set\_cash} 和 \texttt{set\_position}。

本书首先把事件、订单和成交设计为值类型。它们代表某个时刻的事实或意图，复制后仍然具有同样含义，不需要共享身份：

\lstinputlisting[caption={回测项目的基础值类型},firstline=1,lastline=38]{project/include/quant/types.hpp}

值语义让容器、测试和消息传递更直接。若以后测量证明字符串复制是热点，可以引入代码表、驻留字符串或列式布局；在证据出现前，清楚所有权比提前压缩表示更重要。

\section{组合优于继承}

回测引擎拥有撮合器和组合对象，它不需要继承二者。继承适用于真实的可替换子类型关系，并要求基类接口对所有派生类都成立。为了复用几行实现而建立继承层次，往往产生受保护状态、脆弱覆盖和难以理解的生命周期。

运行时多态可以用纯虚接口或类型擦除；编译期多态可以用模板与 concepts。选择依据是替换发生在编译期还是运行时、ABI 是否稳定、代码体积和诊断是否可接受，而不是“现代 C++ 永远不用虚函数”。

\section{撮合与组合的公开契约}

\lstinputlisting[caption={撮合和组合的最小公开接口}]{project/include/quant/execution.hpp}

撮合器拒绝代码不一致、非正数量、流动性不足或非法价格的订单。组合只接受已经形成的成交，并通过快照暴露可观察状态。测试不读取内部散列表，因而未来可以把持仓替换为排序向量或专用簿记结构而不改行为测试。

当前简化组合的权益只覆盖请求快照的单一代码。这是教学边界，不是暗藏假设：第 13 章会明确其适用范围和扩展到多资产估值所需的价格映射、币种与缺失行情政策。

\section{构造后即有效}

理想对象在构造完成后就满足不变量。若配置可能失败，可使用工厂返回值或错误，而不是创建“半有效对象”再要求调用 \texttt{init()}。对于整数数量，规定正数与方向分离；对于价格，拒绝非有限值；对于初始现金，生产接口还应规定是否允许负数和币种。

\section{接口深度}

深接口用少量调用隐藏大量正确性工作。例如 \texttt{apply\_fill} 隐藏买卖方向、现金符号和持仓更新的一致提交。浅接口只是改名转发，迫使调用者理解内部步骤。代码审查时可以问：调用者要读多少实现，才能安全使用这个类？

\begin{interviewcheck}
为什么事件适合值语义？组合为何只接收成交而不接收订单？何时运行时多态优于模板？“构造后即有效”如何改变错误处理接口？
\end{interviewcheck}

\section{练习}

\begin{enumerate}
  \item 为订单增加唯一标识，决定它属于值的一部分还是外部追踪元数据。
  \item 扩展撮合器支持部分成交，写出剩余数量与事件流动性的契约。
  \item 设计多资产组合快照接口，说明缺失最新价格时选择拒绝、沿用旧价还是标记结果不完整。
  \item 找出一个“Manager”类，把它拆成两个各自拥有明确不变量的组件。
\end{enumerate}

\section{本章小结}

可靠类从职责和不变量开始。事件、订单和成交采用值语义，撮合与组合通过小接口协作，内部数据结构保持可替换。下一章用模板与 concepts 把接口要求变成编译期契约。
